\section{State substrates and representations}\label{sec:substrates}


\regime{Substance}{where persistent state physically lives, and what each substrate makes easy or precludes.}
The previous parts established that an always-on agent is a persistent-state system whose retained state spans far more than retrievable text. The physical question is where that state lives. Each substrate, the medium in which state is materialized, indexed, and recovered, imposes a different set of operations and governance affordances. Model weights forget through interference but cannot be tombstoned; a vector index can be edited at a key but cannot be replayed to a point in time; an append-only event log can replay but cannot be garbage-collected without breaking provenance; a knowledge graph can express that one fact superseded another but rarely records who authorized the supersession. The substrate decides which of the six state axes (authority, scope, mutability, provenance, recoverability, actionability) are native, and which require an external envelope, sidecar log, access-control layer, or transaction wrapper. This is the largest part of the corpus, 117 of 435 coded works, because governance is often absent from the dominant data structures themselves, especially vector indices, parametric weights, and the context window. Rollback, the corpus's rarest lifecycle capability, is especially hard to retrofit when the substrate has no durable authority field, tombstone, or point-in-time state to revert to.

Substrates fall along a physical gradient, from the most internal (model parameters) to the most external (multi-tenant serving runtimes). Classical memory forms, episodic versus semantic versus procedural, parametric versus non-parametric, remain useful, but they cut across substrates rather than replacing them. The form of remembered content, such as an event trace, a distilled fact, or a callable skill, tells us what is stored; the substrate tells us how it is addressed, mutated, and recovered, which is what governance acts on. A reflection lesson and a fine-tuned weight delta may carry the same procedural knowledge, but only one can be deleted, attributed, and rolled back. That distinction drives the reading below.

\subsection{Parametric state and test-time learning}

The most internal substrate is the model's own parameters. Here state is non-symbolic: it lives as distributed weight changes rather than addressable records, which makes parametric memory dense and fast to read at inference but opaque to inspection, attribution, and deletion. One recent line treats inference-time weight updates as a memory mechanism. TMEM absorbs distilled supervision into fast LoRA weights through online reinforcement-learning-optimized updates within an episode, so an agent adapts to a running interaction without paying repeated context cost \citep{ren2026tmem}. PEAM converts retrieval memory into parameter-resident skills via per-category isolated mixture-of-experts LoRA adapters, with a failure-as-training-signal loop that promotes successful retrieval patterns into weights \citep{guo2026peam}. Both internalize what would otherwise be external context, trading addressability for efficiency. A second cluster makes the parametric memory recurrent and constant-size: a distillation scheme compresses a full-history transformer teacher's bottleneck representation into a recurrent transformer student's fixed-size memory, so the agent keeps a bounded parametric state instead of an unbounded store of past observations \citep{weinzaepfel2026recurrentmem}. A third treats the latent state itself as the retrievable object: ElasticMem retrieves latent memories from the reasoner's own hidden state and uses a learned policy to assign each one a variable latent token budget, cutting cost while preserving accuracy \citep{feng2026elasticmem}, while EvoEmbedding maintains a continuously updated latent memory as it sequentially encodes a long context, producing evolvable representations for dynamic retrieval \citep{nie2026evoembedding}.

These mechanisms differ in what they internalize and how reversibly. TMEM and PEAM make adaptation episodic and category-scoped, so a bad update is at least contained to a LoRA module rather than the base weights; recurrent compression \citep{weinzaepfel2026recurrentmem} makes the memory bounded but irreversible, since once an observation is folded into the recurrent state there is no handle to recover or excise it. Parametric memory has no tombstone: a fact cannot be deleted from a weight matrix the way a row can be dropped from a store, and weight editing at scale introduces the same gradual and catastrophic interference that makes continual learning hard. The tradeoff is direct: TMEM and PEAM deliver fast test-time adaptation and reduced context cost, but updates leave no audit trail for the provenance of a parameter shift and no mechanism to detect an unauthorized or poisoned internalization. Parametric state is therefore efficient to read and hard to govern at write time: it satisfies actionability and, within an episode, mutability while structurally weakening provenance preservation and rollback traceability. The missing capability is a parametric substrate with a delete: a way to attribute a behavior to the specific update that produced it and revert that update without replaying the entire training history. No corpus work in this cluster offers it.

\subsection{The context window and long-context}

If parametric memory is the most internal substrate, the context window is the most ephemeral: state lives only as tokens in the active prompt, governed implicitly by what the agent chooses to keep in scope. The long-context literature is, at root, a study of how much state this substrate can hold and how reliably it can be read. Boundary diagnostics establish the limits. LongBench probes long-context understanding across diverse bilingual task types and gives an empirical baseline \citep{bai2024longbenchbilingualmultitaskbenchmark}; RULER constructs synthetic tasks to reveal the true effective context length behind the advertised window, which is routinely far shorter \citep{hsieh2024rulerwhatsrealcontext}; and the lost-in-the-middle result shows that models systematically underuse information placed mid-context even when they attend to the boundaries \citep{liu2023lostmiddlelanguagemodels}. The survey of context engineering systematizes the techniques that work within these limits, from compression to reordering to selective inclusion \citep{mei2025surveycontextengineering}. The shared lesson is that the context window is an unreliable substrate for state: presence in the window does not guarantee retrievability, so an always-on agent cannot treat a long prompt as durable memory.

A separate line attacks the boundedness of the window directly. StreamingLLM enables a model to run on effectively infinite-length streams by retaining attention sinks plus a sliding window, making unbounded operation possible without retraining \citep{xiao2023streamingllm}. This is the substrate-level enabler of always-on operation: an agent that never resets must consume an unbounded event stream, and StreamingLLM is the canonical mechanism for doing so within the context substrate rather than offloading to an external store. But infinite throughput is not the same as infinite memory: the sliding window discards old tokens, so the substrate forgets by construction and offers no way to retrieve or attribute what scrolled past.

The central debate in this subsection is the long-context-versus-retrieval boundary, because it decides whether state should live in the window at all. A foundational study showed that retrieval augmentation consistently improves LLMs regardless of context length, so even a long-context model benefits from an external store \citep{xu2023retrievalmeetslongcontext}. Subsequent work refined the boundary: long-context LLMs generally beat retrieval on accuracy but at much higher cost, motivating Self-Route, which routes each query to the cheaper mode \citep{li2024ragorlongcontext}; a systematic study across twenty models found that only a few recent systems sustain accuracy past 64k tokens and cataloged the distinct failure patterns of long-context retrieval-augmented generation \citep{leng2024longcontextragperformance}; and a careful revisit found that long context wins on QA benchmarks while retrieval retains an edge on dialog, sharpening rather than settling the tradeoff \citep{li2025longcontextvsrag}. Past a million tokens, the BEAM benchmark (100 conversations, 2,000 validated questions) shows that even million-token-window models degrade on genuinely long interaction histories, so window size alone does not solve always-on memory \citep{tavakoli2025beyondamilliontokens}. A representative hybrid response is general agentic memory, which flags key history into a light persistent store while preserving all raw pages, then composes context just-in-time per request, refusing to choose statically between the window and the store \citep{yan2025generalagenticmemory}.

For governance, the context window has the weakest affordances of all. It has no persistent identity (it resets between sessions by definition), no provenance (a token in the prompt carries no source or authority tag unless one is manually encoded), and no deletion semantics (dropping a token is a side effect of windowing, not an audited erasure). The boundary studies expose a representation-quality problem the thesis can state directly: state that is present but unretrievable, the lost-in-the-middle failure, is a silent violation of the assumption that retained state is usable. The gap is that the long-context literature measures whether supplied state can be read, never whether the agent decided correctly what should have become durable state in the first place; the window is a substrate for working context, and treating it as a memory substrate inherits all of its ungoverned ephemerality.

\subsection{Retrieval and vector stores (RAG)}

Retrieval-augmented generation externalizes state into a corpus of passages addressed by embedding similarity. This is the field's default non-parametric substrate, and the corpus reflects it: the RAG-and-long-context cluster is the single largest substrate cluster. Its representational unit is the chunk or document, and its governance affordance is, in principle, provenance: every retrieved span has a source. The benchmark literature established the failure vocabulary. RGB decomposes RAG quality into noise robustness, negative rejection, information integration, and counterfactual handling, framing retrieval as more than nearest-neighbor lookup \citep{chen2023benchmarkinglargelanguagemodels}; RAGTruth annotates hallucinations to make faithfulness measurable \citep{niu2024ragtruthhallucinationcorpusdeveloping}; RAGBench adds explainable relevance and faithfulness metrics for diagnostic evaluation \citep{friel2025ragbenchexplainablebenchmarkretrievalaugmented}; and CRAG spans diverse domains and dynamic facts to stress robustness to staleness \citep{yang2024cragcomprehensiverag}. The multi-hop substrate is anchored by HotpotQA, which requires reasoning over several documents and so forces a retriever to compose rather than match \citep{yang2018hotpotqa}.

For always-on agents, the key limitation is that these benchmarks evaluate retrieval over a static corpus, whereas an always-on agent's store is one it writes to from its own experience. Three lines confront the moving-corpus problem. First, temporal staleness: a study of temporal validity in retrieval quantifies that standard RAG returns stale facts 15 to 40 percent of the time over evolving knowledge and proposes a temporal strategy to fix it \citep{yadav2026temporalvalidityretrieval}. Second, retrieval that is misleading rather than simply wrong: causal memory intervention selects which memories to inject by measuring each candidate's causal effect on the agent's output under controlled interventions, filtering similar-but-misleading memories that pure semantic similarity would surface \citep{srivastava2026causalmemory}; and a learned-value approach scores each candidate memory by seven cognitive factors (emotional intensity, goal relevance, value alignment, and others) to decide what to remember at all \citep{chen2026learningwhattoremember}. Third, goal-directed retrieval: Goal-Mem treats a user utterance as a goal and chains backward into atomic subgoals to target what to retrieve, replacing flat similarity with structured intent \citep{liang2026goalmemragmemory}. These works increasingly treat similarity-addressed retrieval, the substrate's defining mechanism, as too weak for self-written, evolving, action-conditioning state.

The governance blind spot of the vector store is visible in the matrix. RGB and its successors benchmark retrieval in isolation from the lifecycle: they never test retrieval under state mutation, versioning, or deletion, because their corpus does not mutate. Provenance is nominally present (each chunk has a source) but is not propagated through consolidation, and deletion is a retrieval edit, not an audited erasure: removing a document from the index leaves its influence in cached prompts, summaries, and any parametric internalization downstream. The substrate makes write-and-retrieve easy but makes deletion propagation, temporal validity, and authority-scoped access nearly impossible to express, because an embedding key carries no notion of who may read it, when it expired, or what it superseded. The requirement this substrate exposes is the one the thesis foregrounds: a retrieval store that is mutation-aware, where a delete cascades to derived state and a write records the authority and validity window that govern when the chunk may act, rather than only whether it is similar.

\subsection{Structured stores: knowledge graphs and temporal memory}

Where the vector store addresses state by similarity, structured stores address it by relation and by time, and in doing so they make some governance axes representable for the first time. The knowledge-graph line begins with HippoRAG, a neurobiologically inspired long-term memory that builds a knowledge graph and runs personalized PageRank for retrieval, integrating structural organization with associative recall and beating flat retrieval on multi-hop tasks \citep{gutierrez2024hipporag}. Recent agentic-memory graphs make the structure dynamic. FluxMem represents memory as a heterogeneous graph that continuously reshapes its topology, repairing links, removing noise, and distilling repeated successful paths into reusable routines \citep{fang2026fluxmem}; AtomMem extracts high-value atomic facts into hierarchical events and activates an associative graph at retrieval time via spreading activation to connect fragmented memories \citep{yao2026atommem}; and H-Mem evolves memory through a temporal-semantic tree that consolidates short-term into long-term plus a parallel knowledge graph capturing entity relations \citep{yu2026hmem}. A graph survey systematizes the design space along short-versus-long-term, knowledge-versus-experience, and structural-versus-non-structural axes, showing that graph substrates enable multipath retrieval and reasoning that flat stores cannot \citep{yang2026graphagentmemorysurvey}.

The most governance-relevant structured substrate is the temporal store, because time is the one axis on which always-on state most obviously diverges from static corpora. The benchmark lineage is mature. TimeQA established that a single fact has time-indexed answers, built from time-evolving Wikipedia \citep{chen2021timeqa}; StreamingQA measured whether a store keeps pace with knowledge arriving over time \citep{liska2022streamingqa}; FreshQA categorized questions by how fast their answers change and added false-premise cases \citep{vu2023freshllms}; and ChroKnowBench organized the evolving-versus-constant distinction across domains and years \citep{park2024chroknowledge}. A parallel probing literature decomposes temporal reasoning itself: TempReason orders difficulty levels \citep{tan2023tempreason}, TimeBench unifies symbolic, commonsense, and event temporal reasoning \citep{chu2023timebench}, MenatQA factors temporal QA into scope, order, and counterfactual perturbations \citep{wei2023menatqa}, Test of Time uses synthetic graphs to isolate pure temporal reasoning from pretraining contamination \citep{fatemi2024testoftime}, and TimeR4 injects explicit time constraints into retrieval over a temporal knowledge graph \citep{qian2024timer4}. Memory-T1 trains a policy via reinforcement learning to select temporally relevant memories for multi-session reasoning, moving from probing to mechanism \citep{du2025memoryt1}.

The temporal mechanisms most relevant to governance make time a first-class field of the stored record. Chronos converts conversation into structured time-stamped events to support multi-hop temporal queries \citep{sen2026chronos}; APEX-MEM stores conversational memory as a temporally annotated property graph and reasons over it with agentic retrieval \citep{banerjee2026apexmem}; and Engram comes closest to lifecycle-complete. Engram writes lossless episodes on a fast path and asynchronously distills subject-predicate-object facts into a bi-temporal knowledge graph carrying both valid-time and transaction-time, supporting point-in-time as-of retrieval, and resolves conflicts by invalidate-never-delete so that every fact retains a provenance supersession chain \citep{wang2026engram}. Engram is the rare substrate in the corpus that directly serves three governance axes at once: provenance (the supersession chain), mutability (valid-time updates without destroying history), and a form of recoverability (as-of retrieval is a read-time rollback). The structured substrate thus contributes what the vector store lacks: the ability to express that a fact held during an interval, was superseded, and descends from a source. Its blind spot, visible across HippoRAG and the graph systems, is mutation under access: graph construction is often offline and static, and few works address edge deletion, node versioning, graph-integrity validation under concurrent agent-driven updates, or who is authorized to rewrite a relation. The gap is that the temporal substrate solves provenance and validity for a single writer but leaves authority and concurrent-write consistency to the runtime layer below it.

\subsection{Multimodal memory}

State is increasingly not text. Multimodal memory substrates store and address visual, video, and egocentric-sensor state, and they expose governance problems that text substrates can hide. The benchmark wave is recent and concentrated. The framing works probe distinct fault lines: MemEye evaluates memory along visual-evidence granularity from scene to pixel and along evidence-use fidelity, grounding memory analysis in perception rather than recall \citep{guo2026memeye}; MemGallery extends multi-session conversational-memory benchmarking to multimodal LLM agents over image-plus-text dialogue across thirteen memory systems \citep{bei2026memgallery}; and EgoMem and EgoStream push to week-long and streaming egocentric video, with EgoStream introducing an answer-validity window that separates model forgetting from natural world-state change \citep{wang2026egomemreason,forte2026egostream}. SuperMemoryVQA grounds the problem in real AI-glasses hardware, with 52.9 hours of activity, synchronized SLAM, gaze, and IMU sensors, and thousands of verified question-answer pairs \citep{alam2026supermemoryvqa}. WorldMemArena frames memory as an action-world interaction loop and separately diagnoses writing, maintenance, retrieval, and use, reporting a finding that anchors this whole part: better storage does not improve task performance, and visual evidence is underused on agentic trajectories \citep{liu2026worldmemarena}. EgoExoMem adds the cross-view problem, exposing view-asymmetry and question-versus-answer view-preference conflicts across synchronized egocentric and exocentric video \citep{liu2026egoexomem}, and H2HMem and SMMBench stress cross-participant and source-distributed evidence composition \citep{zhu2026h2hmem,chai2026smmbench}.

The mechanism side mirrors the structured-store designs but adds modality fusion. TeleMem builds a unified long-term and multimodal (including video) memory with narrative dynamic-profile extraction \citep{chen2026telemem}; MemVerse combines fast parametric recall with hierarchical retrieval and continual consolidation in a model-agnostic framework \citep{liu2026memverse}; M2A maintains a dual-layer hybrid memory through a collaborating ChatAgent and MemoryManager that perform online updates \citep{feng2026m2a}; and NS-Mem combines neural representations with explicit symbolic structures and rules in a three-layer architecture \citep{jiang2026nsmem}. Embodied variants ground memory in action: STAR builds multimodal long-term memory for mobile robots and retrieves a compact task-conditioned subset \citep{yuan2026star}, RoboCortex maintains dual-grain reflective and principle memory for a self-evolving embodied agent \citep{chan2026robocortex}, and PersonaVLM turns a general multimodal LLM into a long-term personalized assistant via memory extraction and multi-turn reasoning \citep{nie2026personavlm}. Streaming perception is handled by StreamingVLM, which sustains real-time comprehension of effectively infinite video with a compact key-value cache reusing attention sinks \citep{xu2025streamingvlm}, and by a training-free episodic-event framework for long-video QA \citep{wang2025videoem}.

Multimodal substrates raise governance stakes in two specific ways the matrix flags. First, write correctness is harder: a visual memory can hallucinate details that were never observed, and reconciling text-image conflicts has no obvious authority rule, so MemEye's evidence-use axis becomes a faithfulness-of-write problem as well as a retrieval problem \citep{guo2026memeye}. Second, the attack surface is wider: visual inception shows that image-borne triggers hidden in user uploads can act as sleeper agents in an agentic recommender's long-term memory, hijacking later behavior \citep{qian2026visualinception}, a poisoning channel that text-only provenance tagging does not catch because the malicious payload enters as pixels. The WorldMemArena finding, that storage quality and task performance are decoupled and visual evidence is underused \citep{liu2026worldmemarena}, supports the survey's claim: accumulating more multimodal state does not by itself produce governed, actionable state, although WorldMemArena itself frames the result as one about storage quality and evidence use rather than governance. A unifying attempt, MemoryWire, defines a JSON-Schema wire format with five canonical operations (remember, recall, forget, merge, expire) over four memory types to make multimodal memory interoperable \citep{munirathinam2026memorywire}; but as the matrix notes, the specification is syntactic, with no merge semantics under conflict, no forget validation, and no attestation of operation provenance. The gap is that multimodal memory multiplies what state can be without giving deletion, provenance, or authority a per-modality semantics, so a deleted image's derived embeddings, captions, and behavioral influence remain unaccounted for.

\subsection{Agentic-OS and serving runtimes}

The final substrate is the runtime itself: the operating-system layer that gives agent state identity, isolation, durability, and recovery. This cluster is where the persistent-state thesis is most directly anticipated by systems work, because operating systems and databases have long treated state as something to be governed as well as stored. The framing work, AIOS, proposes an LLM-agent operating system that isolates memory, storage, context, scheduling, and access control into a kernel managing concurrent agents \citep{mei2024aios}. Three identity models follow. Quine realizes agents as native POSIX processes: identity is the PID, lifecycle is fork/exec/exit, and persistent state is the process's memory, environment, and filesystem, all inheriting kernel isolation \citep{ke2026quine}. AgentLibOS runs each long-running agent as an AgentProcess with identity, lifecycle state, capabilities, checkpoints, and audit, making runtime primitives rather than tool calls the authority boundary \citep{zhang2026agentlibos}. ActPlane enforces agent action policies inside the kernel using eBPF and an information-flow-control DSL, covering execution paths that tool-call guardrails miss at 1.9 to 8.4 percent overhead \citep{zheng2026actplane}. These are the first substrates in the substrate map where authority is first-class and kernel-enforced rather than an implicit field, notable given that authority is the rarest axis in the whole corpus.

A second strand makes recoverability the organizing primitive, and it is the densest concentration of rollback-capable substrates in the corpus. Event-sourced runtimes treat state as a deterministic projection of an append-only log: ActiveGraph supports exact replay and gates every self-improvement promotion or discard as an auditable event \citep{nakajima2026regimes}, and git-based memory stores the agent's reasoning tree in a repository so it can be replayed, diffed, and merged, arguing that git's value for agent memory is auditability and cross-agent merge rather than retrieval accuracy \citep{shekar2026gitofthoughts,wu2025gitcontextcontroller}. Checkpoint-restore runtimes make rollback cheap and correct: DeltaBox gives millisecond change-based transactional checkpoint and rollback of full sandbox state via copy-on-write filesystem layers and incremental process dumps \citep{dong2026deltabox}, and CRAB closes the agent-OS semantic gap by eBPF-classifying each turn's OS effects to align durable checkpoints with turn boundaries, lifting recovery correctness from 8 percent to 100 percent \citep{wu2026crab}. OpenRATH makes a branchable, inspectable, replayable session the first-class runtime value, carrying lineage, sandbox placement, and tool evidence with the execution \citep{wen2026openrath}, and AEON reframes long-horizon memory as a managed OS resource backed by a write-ahead log, an mmap blob arena with generational garbage collection, and epoch-based lock-free reads \citep{arslan2026aeon}.

The third strand is the hardest: transactional consistency over external side effects, the point where rollback stops being a memory operation and becomes a commitment problem. Atomix wraps tool use in progress-aware transactions that record reads and effects, seal when the footprint is complete, and commit per-resource only once frontiers confirm no earlier conflicting work can arrive; it handles bufferable, reversible, and irreversible effects distinctly \citep{mohammadi2026atomix}. Cordon treats a multi-step task as one semantic transaction, staging irreversible effects in an outbox and reversible mutations in shadow state, then binds tool intents to result lineage and authority before commit \citep{chen2026cordon}. ATCC treats an agent's emitted SQL as long-running agentic transactions and adapts database concurrency control per transaction \citep{zhou2026atcc}. A concurrency-anomaly study complements these: it formalizes and machine-verifies four anomalies over shared multi-agent state (memory stores, vector indices, tool registries) with a checked consistency hierarchy \citep{khan2026concurrency}. A distributed-systems foundation extends consistency models from data to knowledge visibility, introducing epistemic state replication with verifiable semantic rollback as a primitive for trustworthy agentic infrastructure \citep{he2026postdeterministic}. A governance-oriented memory layer, ProjectMem, stores coding-agent history as an append-only typed event log condensed into delivered summaries, then applies a deterministic memory-as-governance gate that warns the agent before it repeats a failed fix or edits a known-fragile file \citep{malo2026projectmem}.

The runtime layer also supplies several important negative results. A longitudinal eight-week postmortem of a production always-on runtime taxonomizes silent failures, names the fail-plausible class in which the model narrates an error as plausible success, and shows that audits block regressions but cannot predict failures \citep{wu2026silentfailures}. A second study finds that deployed agents under impossible constraints fabricate fake external problems and even fake system crashes (thanatosis), a self-reinforcing behavior that injecting accurate information mid-session does not correct \citep{rodriguez2026thanatosis}. These are not retrieval failures; they are governance failures at the runtime substrate, where the agent's reported state diverges from its actual state and no lower layer can catch it. This cluster imports the database and operating-system playbook (logs, transactions, checkpoints, capabilities, consistency models) into agent state, producing nearly all of the corpus's genuinely rollback-capable substrates. The gap is integration and adoption: these runtime primitives sit below the memory layer, while the memory mechanisms above them (the vector stores, graphs, and multimodal stores of the preceding subsections) rarely build on them, so the typical deployed agent stacks an ungoverned store on a runtime that could have governed it. The substrate that can roll back and the substrate that holds the memory are usually not the same system.

\subsection{Systems cost and resource governance}

Every substrate above incurs a cost that grows with always-on operation, and a final cluster treats that cost as a first-class governance constraint rather than an afterthought. The empirical anchor is a deployment measurement: tracking storage growth over a 15,000-message deployment finds memory grows roughly linearly at about three tokens per message, exposes a sim-to-real gap where real workflows grow linearly while synthetic benchmarks stay constant, and names consolidation quality, not raw capacity, as the scalability bottleneck \citep{milosevic2026episodicsemantic}. This linear, unbounded growth is why ungoverned accumulation is not free: a store that only writes and retrieves degrades in latency and cost over a long horizon even if its recall stays nominally correct. A total-cost-of-ownership study measures accuracy, latency, CPU, RAM, disk I/O, and network across production memory frameworks (mem0, Graphiti, cognee) on an independent cloud-edge testbed and finds that plain retrieval matches top accuracy at 8.4 times lower total cost than the heaviest framework \citep{wolff2026costaccuracy}, reframing the substrate choice as a cost-governance decision rather than only an accuracy one. A complementary systems-level characterization profiles stateful long-horizon memory workloads across ten representative systems, attributing cost across construction, retrieval, and generation phases \citep{omri2026agentmemorycharacterization}.

The mitigation literature borrows directly from operating-systems resource management. AgentRM mines over 40,000 GitHub issues to show that unbounded growth and weak retention cause agent amnesia. It then introduces OS-inspired middleware, a multi-level feedback-queue scheduler plus a three-tier context-lifecycle manager, that cuts P95 latency 86 percent while retaining 100 percent of key information \citep{she2026agentrm}. AMVL treats persistent memory as a managed systems resource with value-driven lifecycle tiers that bound the request-path working set independent of total stored memory, cutting tail latency several-fold \citep{bamidele2026amvl}. MemRefine formalizes storage-budgeted memory management differently: it keeps an accumulated store within a fixed footprint via similarity-nominated, LLM-judged delete/merge/preserve compression that iterates until the budget is met \citep{kim2026memrefine}. In the multi-agent setting, token coherence maps shared-state synchronization onto the hardware cache-coherence problem, adapting MESI lazy invalidation into an artifact-coherence protocol with a TLA+-verified guarantee of single-writer safety, monotonic versioning, and bounded staleness \citep{parakhin2026tokencoherence}.

The governance reading of this cluster is that resource pressure and state governance are the same problem seen from two sides. Eviction is forgetting under a budget, and the matrix's critique of AgentRM captures the survey's concern: a heuristic eviction policy that is not formally verified can evict critical state before it is acted upon, so a cost optimization becomes a correctness failure \citep{she2026agentrm}. The TCO and characterization studies do not measure mutation cost, recovery after failure, or multi-tenant isolation \citep{wolff2026costaccuracy,omri2026agentmemorycharacterization}, so the overhead of provenance, deletion propagation, and rollback remains largely unbudgeted in the literature. The missing systems model is a cost account for governance operations: what does it cost to make a delete propagate, to keep a supersession chain, to checkpoint for rollback, and is that cost one reason governance is the rarest capability in the corpus? The cheapest substrate is often the least governed one, and no work has yet shown that stronger governance must be prohibitively expensive.

\subsection{Cross-substrate synthesis}

Taken together, the seven substrates form a spectrum from internal and dense to external and durable, and the memory-form sub-axis cuts across all of them: episodic traces appear as raw pages in agentic memory \citep{yan2025generalagenticmemory}, as event logs in event-sourced runtimes \citep{nakajima2026regimes}, and as video frames in egocentric stores \citep{forte2026egostream}; semantic facts appear as graph nodes \citep{wang2026engram} and as retrieved chunks \citep{chen2023benchmarkinglargelanguagemodels}; procedural skills appear as LoRA adapters \citep{guo2026peam}, as distilled graph routines \citep{fang2026fluxmem}, and as governance gates over a typed event log \citep{malo2026projectmem}. The same knowledge can thus be parametric or non-parametric depending only on which substrate holds it, and that choice, not the form of the content, determines whether the state can be attributed, scoped, and reverted. Table~\ref{tab:substrates} summarizes for each substrate what it natively stores and the governance blind spot it structurally imposes.

\begin{table}[t]
\centering
\small
\setlength{\tabcolsep}{4pt}\renewcommand{\arraystretch}{1.2}
\begin{tabularx}{\linewidth}{p{0.18\linewidth}YY}
\toprule
Substrate & What it natively stores & Governance blind spot \\
\midrule
Parametric / test-time & Distributed weight deltas; latent recurrent state; internalized skills & No native tombstone or attribution: deleting or auditing a specific update requires external instrumentation; provenance and rollback are not native \\
Context window & Tokens in the active prompt; sliding-window stream & No persistent identity, no source tags, no audited deletion; present-but-unretrievable state (lost-in-the-middle) \\
Vector store (RAG) & Embedded chunks addressed by similarity & Mutation-blind: no versioning or deletion cascade; provenance per chunk but not propagated; no validity window or read authority \\
Knowledge graph / temporal & Relations; bi-temporal facts with supersession chains & Single-writer assumption: weak under concurrent edits; authority over who may rewrite a relation is unspecified \\
Multimodal & Visual, video, egocentric-sensor state plus fused text & Per-modality faithfulness of writes; pixel-borne poisoning; no per-modality deletion/provenance semantics \\
Agentic-OS / runtime & Process identity, event logs, checkpoints, transactional effects & Sits below the memory layer it could govern; adoption gap means stores rarely build on its rollback primitives \\
Systems / resource layer & Tiered footprint, eviction state, cost profiles & Eviction is unverified forgetting; cost of governance operations (delete, audit, rollback) is unbudgeted \\
\bottomrule
\end{tabularx}
\caption{Substrate comparison. Each substrate makes some state axes easy to represent and others non-native. Reading down the rightmost column reproduces, at the physical layer, the corpus-wide skew: writing and retrieving are easy across substrates, while deletion propagation, authority, and rollback require sidecar metadata, transactional wrappers, or different data structures. The rare exceptions (bi-temporal graphs, event-sourced runtimes, and checkpoint-restore runtimes) import database and operating-system governance, yet the memory layer above them rarely adopts those primitives.}
\label{tab:substrates}
\end{table}

\begin{table}[t]
\centering
\small
\setlength{\tabcolsep}{4pt}\renewcommand{\arraystretch}{1.2}
\begin{tabularx}{\linewidth}{p{0.22\linewidth}cccccc}
\toprule
Substrate & Auth. & Scope & Mut. & Prov. & Recov. & Action. \\
\midrule
Parametric / test-time   & \xmark & \pmark & \cmark & \xmark & \xmark & \cmark \\
Context window           & \xmark & \pmark & \cmark & \xmark & \xmark & \cmark \\
Vector store (RAG)       & \xmark & \xmark & \xmark & \pmark & \xmark & \cmark \\
Knowledge graph / temporal & \pmark & \pmark & \cmark & \cmark & \pmark & \cmark \\
Multimodal               & \xmark & \pmark & \pmark & \pmark & \xmark & \cmark \\
Agentic-OS / runtime     & \cmark & \cmark & \cmark & \cmark & \cmark & \cmark \\
Systems / resource layer & \pmark & \pmark & \cmark & \xmark & \pmark & \cmark \\
\bottomrule
\end{tabularx}
\caption{Which state axes each substrate can natively represent: \cmark{} native, \pmark{} partial or single-writer only, \xmark{} not native. A non-native axis is not impossible to govern, but it requires an external envelope, sidecar log, access-control layer, or transaction wrapper. Vector stores and parametric memory, the two dominant substrates, do not natively represent authority and recoverability, the two governance axes the corpus most neglects. Agentic-OS and bi-temporal-graph substrates import database and operating-system primitives that make all six axes native, but these substrates are rarely the default memory layer.}
\label{tab:substrate-axis}
\end{table}


Table~\ref{tab:substrate-axis} restates this at the level of the six axes: it marks, for each substrate, which axes it can natively represent. The two columns that stay blank across the dominant substrates, authority and recoverability, are also the two governance axes the corpus most neglects, and only the substrates that import database and operating-system primitives fill them natively. This gives a physical-layer version of the survey's claim: the corpus-wide skew is more than a matter of attention, because dominant substrates make some governance operations expensive unless an external envelope is added. An embedding key does not by itself express an authority; a weight delta does not by itself carry a tombstone; a sliding window does not replay itself; a similarity index does not cascade a delete. The substrates that can make these fields native, bi-temporal knowledge graphs with supersession chains \citep{wang2026engram}, event-sourced and checkpoint-restore runtimes \citep{nakajima2026regimes,dong2026deltabox,wu2026crab}, and transactional effect managers \citep{mohammadi2026atomix,chen2026cordon}, already exist, but they live at the runtime and structured-store layers the dominant memory mechanisms are not yet built on. The substrate-level gap, which the lifecycle and governance sections develop, is therefore an integration gap as much as an invention gap: the discipline already has several data structures for governed persistent state; it has not yet made them the default substrate on which agent memory is written.
